% Kandydat restrukturyzacji; nie jest aktywnym rozdziałem Atlasu.
\section{Interpolacja zespolona i nierówność Hausdorffa--Younga}
\providecommand{\atlasref}[2]{#2}

\subsection{Interpolacja zespolona i nierówność Hausdorffa--Younga}

Wynik dla $L^1$ daje kontrolę normy supremum transformaty,
a twierdzenie Plancherela zachowanie normy $L^2$.
Pomiędzy tymi przypadkami można uzyskać oszacowania w
$L^{p'}$, gdzie $p'$ jest wykładnikiem sprzężonym do $p$.
Potrzebujemy jednak interpolacji norm operatorów, nie tylko
interpolacji norm jednej funkcji.

\begin{lemat}[Lemat o trzech prostych]
\label{lem:dt-three-lines}
Niech $F$ będzie ograniczoną funkcją ciągłą na domkniętym
pasie $0\le\operatorname{Re}z\le1$, holomorficzną wewnątrz
pasa. Załóżmy, że dla każdego $t\in\R$ zachodzą
$|F(it)|\le M_0$ oraz $|F(1+it)|\le M_1$.
Dla $0<\theta<1$ mamy
\[
 |F(\theta)|\le M_0^{1-\theta}M_1^\theta.
\]
\end{lemat}
\begin{proof}
Najpierw zakładamy $M_0,M_1>0$. Funkcja
$G(z)=F(z)\exp(-(1-z)\log M_0-z\log M_1)$
jest ograniczona na pasie i ma moduł co najwyżej jeden
na jego pionowych brzegach. Dla $\varepsilon>0$ połóżmy
$H_\varepsilon(z)=G(z)\exp(\varepsilon(z^2-1))$.
Jeżeli $z=x+iy$, $0\le x\le1$, to moduł dodatkowego
czynnika wynosi $\exp(\varepsilon(x^2-y^2-1))$.
Na pionowych brzegach jest więc nie większy od jedności.
Na poziomych odcinkach $y=\pm R$ moduł $H_\varepsilon$
nie przekracza $\|G\|_\infty e^{-\varepsilon R^2}$,
co dla dużych $R$ jest mniejsze od jedności.
Z zasady maksimum modułu na prostokącie otrzymujemy
$|H_\varepsilon(\theta)|\le1$.
Przejście do granicy $\varepsilon\downarrow0$ daje
$|G(\theta)|\le1$. Dla $M_j=0$ stosujemy udowodniony
wynik ze stałymi $M_j+\delta$ i przechodzimy do
granicy $\delta\downarrow0$.
\end{proof}

\begin{twierdzenie}[Twierdzenie Riesza--Thorina]
\label{thm:dt-riesz-thorin}
Niech $(X,\mu)$ i $(Y,\nu)$ będą przestrzeniami
z miarami $\sigma$-skończonymi. Niech $T$ będzie
operatorem liniowym określonym na zespolonych funkcjach
prostych o nośniku skończonej miary w $X$, o wartościach
w klasach funkcji mierzalnych w $Y$.
Załóżmy, że dla $j=0,1$ oraz wszystkich takich $f$
\[
 \|Tf\|_{q_j}\le M_j\|f\|_{p_j},
 \qquad 1\le p_j<\infty,\quad 1\le q_j\le\infty.
\]
Dla $0<\theta<1$ definiujemy $p,q$ przez
\[
 \frac1p=\frac{1-\theta}{p_0}+\frac{\theta}{p_1},
 \qquad
 \frac1q=\frac{1-\theta}{q_0}+\frac{\theta}{q_1},
\]
przyjmując $1/\infty=0$. Wówczas $T$ przedłuża się
jednoznacznie do ograniczonego operatora
$L^p(X)\to L^q(Y)$, którego norma nie przekracza
$M_0^{1-\theta}M_1^\theta$.
\end{twierdzenie}
\begin{proof}
Przyjmijmy najpierw dodatnie $M_j$. Wystarczy rozważyć
niezerowe funkcje proste $f,g$ o nośnikach skończonej
miary, unormowane przez $\|f\|_p=\|g\|_{q'}=1$.
Na zbiorach, na których funkcje są niezerowe, piszemy
$f=\omega_f|f|$ oraz $g=\omega_g|g|$, gdzie
$|\omega_f|=|\omega_g|=1$. Dla $z$ w pasie definiujemy
\[
\begin{aligned}
 a(z)&=p\bigl((1-z)/p_0+z/p_1\bigr),\\
 f_z&=\omega_f\exp(a(z)\log|f|).
\end{aligned}
\]
Poza zbiorem $\{f\ne0\}$ przyjmujemy $f_z=0$.
Jeśli $q'<\infty$, definiujemy analogicznie
\[
 b(z)=q'\bigl((1-z)/q_0'+z/q_1'\bigr),
 \qquad
 g_z=\omega_g\exp(b(z)\log|g|),
\]
z wartością zero poza $\{g\ne0\}$.
Jeżeli $q'=\infty$, to $q=q_0=q_1=1$ i bierzemy
po prostu $g_z=g$.
Mamy $f_\theta=f$, $g_\theta=g$. Na brzegach pasa
$\|f_{j+it}\|_{p_j}=1$. Istotnie, podniesienie
modułu do potęgi $p_j$ daje $|f|^p$.
Podobnie $\|g_{j+it}\|_{q_j'}\le1$; jeśli
$q_j'=\infty$ w przypadku $q'<\infty$, moduł
na niezerowym nośniku jest równy jeden.

Rozważmy biliniowe parowanie
\[
 F(z)=\int_Y (Tf_z)(y)g_z(y)\,\mathrm d\nu(y).
\]
Jest ono skończone: rozpisując $f_z,g_z$ na skończone
sumy funkcji charakterystycznych, otrzymujemy całki
$T\chi_A$ po zbiorach skończonej miary, które istnieją
na mocy oszacowań brzegowych i nierówności Höldera.
Liniowość $T$ pokazuje, że $F$ jest skończoną sumą
funkcji wykładniczych o wykładnikach liniowych w $z$.
Ponieważ współczynniki przy $z$ w tych wykładnikach
są rzeczywiste, $F$ jest holomorficzna i ograniczona
na pasie. Na jego brzegach, z nierówności Höldera,
\[
 |F(j+it)|\le
 \|Tf_{j+it}\|_{q_j}\|g_{j+it}\|_{q_j'}\le M_j.
\]
Lemat~\ref{lem:dt-three-lines} daje
$|\int_Y (Tf)g\,\mathrm d\nu|\le M_0^{1-\theta}M_1^\theta$.
Po uwzględnieniu dowolnych norm $f,g$ otrzymujemy
\[
 \left|\int_Y(Tf)g\,\mathrm d\nu\right|
 \le M_0^{1-\theta}M_1^\theta\|f\|_p\|g\|_{q'}.
\]

Wyjaśnijmy, dlaczego wystarcza testowanie funkcjami
prostymi. Niech $u$ będzie funkcją mierzalną skończoną
prawie wszędzie, lokalnie całkowalną na zbiorach
skończonej miary, oraz niech
$|\int ug\,\mathrm d\nu|\le C\|g\|_{q'}$ dla takich
funkcji prostych $g$. Jeżeli $1<q<\infty$, wybieramy
rosnące wyczerpanie $Y$ zbiorami $B_n$ skończonej miary
i kładziemy $E_n=B_n\cap\{|u|\le n\}$. Testujemy
ograniczoną funkcją
$\chi_{E_n}\overline{\omega_u}|u|^{q-1}$, z fazą
$\omega_u=u/|u|$ na $\{u\ne0\}$ i zerem w pozostałych
punktach. Jest ona jednostajną granicą funkcji prostych
o nośnikach w $E_n$, więc nierówność przechodzi do
granicy zarówno w całce, jak i w normie $L^{q'}$.
Dostajemy
\[
 \int_{E_n}|u|^q\,\mathrm d\nu
 \le C\left(\int_{E_n}|u|^q\,\mathrm d\nu\right)^{1/q'}.
\]
Wynika stąd $\|\chi_{E_n}u\|_q\le C$, a twierdzenie
o zbieżności monotonicznej daje $\|u\|_q\le C$.
Dla $q=1$ stosujemy test $\chi_{E_n}\overline{\omega_u}$
o normie supremum nie większej od jedności.
Dla $q=\infty$ i $\eta>0$, gdyby zbiór
$\{|u|>C+\eta\}$ miał dodatnią miarę, miałby podzbiór
$B$ o dodatniej skończonej mierze. Test
$\chi_B\overline{\omega_u}$, znów przybliżony jednostajnie
funkcjami prostymi, prowadziłby do sprzeczności
$(C+\eta)\nu(B)\le C\nu(B)$.
Stosując ten argument do $u=Tf$, otrzymujemy żądane
oszacowanie normy. Przypadek zerowego $M_j$ wynika
z zastąpienia go przez $M_j+\delta$.
Ponieważ $p<\infty$, funkcje proste o nośniku skończonej
miary są gęste w $L^p$, a $L^q$ jest zupełna.
Przedłużenie przez granice jest zatem dobrze określone,
ograniczone i jednoznaczne.
\end{proof}

\begin{twierdzenie}[Nierówność Hausdorffa--Younga]
\label{thm:dt-hausdorff-young}
Niech $1\le p\le2$, a $p'$ będzie wykładnikiem sprzężonym.
Dla unitarnej normalizacji transformacji Fouriera
z Definicji~\ref{def:fourier-schwartz17} zachodzi
\[
 \|\widehat f\|_{p'}
 \le (2\pi)^{-d(1/p-1/2)}\|f\|_p,
 \qquad f\in L^p(\R^d).
\]
Transformata jest jednoznacznym ciągłym przedłużeniem
operatora całkowego z funkcji prostych o nośniku
skończonej miary, zgodnym z transformacją na $\mathcal S'$.
\end{twierdzenie}
\begin{proof}
Dla $p=1$ stosujemy Twierdzenie~\ref{thm:dt-fourier-l1-rl},
dla $p=2$ Twierdzenie~\ref{thm:plancherel22}.
Dla $1<p<2$ korzystamy z Twierdzenia~\ref{thm:dt-riesz-thorin}
z parametrami
\[
 (p_0,q_0,M_0)=(1,\infty,c_d),\qquad
 (p_1,q_1,M_1)=(2,2,1).
\]
Oba oszacowania dotyczą tego samego operatora dzięki
Stwierdzeniu~\ref{prop:dt-fourier-l1-l2-compat}.
Mamy $1/p=1-\theta/2$, $1/q=\theta/2$, zatem
$q=p'$ i $1-\theta=2/p-1$.
Norma interpolowana wynosi
$c_d^{\,2/p-1}=(2\pi)^{-d(1/p-1/2)}$.

Aby sprawdzić zgodność z dystrybucjami, wybieramy
funkcje proste $f_n$ o nośnikach skończonej miary
zbiegające do $f$ w $L^p$. Wówczas
$\widehat f_n\to\widehat f$ w $L^{p'}$.
Dla $\varphi\in\mathcal S$ nierówność Höldera daje
\[
 \left|\int_{\R^d}(f_n-f)\mathcal F\varphi\,
 \mathrm d\lambda_d\right|
 \le\|f_n-f\|_p\|\mathcal F\varphi\|_{p'}\to0,
\]
a analogicznie parowanie $\widehat f_n-\widehat f$
z $\varphi$ dąży do zera. Każde $f_n$ należy do $L^1$,
więc stosujemy
Stwierdzenie~\ref{prop:dt-fourier-l1-l2-compat}
i przechodzimy do granicy w tożsamości biliniowej.
Dostajemy $\mathcal FT_f=T_{\widehat f}$.
\end{proof}

\begin{uwaga}[Zakres i stała]
\label{rem:dt-hy-constant}
Powyższa stała pochodzi z interpolacji dwóch norm
brzegowych. Nie twierdzimy, że jest optymalna dla
$1<p<2$. Ostre oszacowanie Babenki--Becknera
wymaga odrębnego dowodu.
Nie można natomiast rozszerzyć samego oszacowania
$\|\widehat f\|_{p'}\le C\|f\|_p$ na $p>2$
przez formalną zamianę wykładników. Kontrprzykład
do takiego rozszerzenia zostanie omówiony poniżej.
\end{uwaga}

\begin{stwierdzenie}[Brak analogicznego oszacowania dla $p>2$]
\label{prop:dt-hy-fails}
Dla $2<p\le\infty$ nie istnieje stała $C$ taka, że
$\|\widehat f\|_{p'}\le C\|f\|_p$ dla wszystkich
$f\in\mathcal S(\R^d)$.
\end{stwierdzenie}
\begin{proof}
Niech $\psi$ będzie niezerową funkcją gładką
o nośniku w kuli o promieniu $1/4$ i niech
$\phi=\mathcal F^{-1}\psi\in\mathcal S$.
Dla $N\ge1$ i znaków $\varepsilon_j\in\{-1,1\}$
kładziemy
\[
 f_N(x)=\phi(x)\sum_{j=1}^N\varepsilon_j e^{ijx_1}.
\]
Transformaty składników mają parami rozłączne nośniki,
więc dla $q=p'<2$ (także $q=1$)
\[
 \|\widehat f_N\|_q=N^{1/q}\|\psi\|_q.
\]
Dla skończonego $p>2$ wybierzmy znaki niezależnie,
z równymi prawdopodobieństwami obu wartości.
Dla $m$ całkowitego z $2m\ge p$ rozwijamy
$\left|\sum_j\varepsilon_j e^{ijx_1}\right|^{2m}$.
Wartość oczekiwana składnika znika, jeśli którykolwiek
indeks występuje nieparzystą liczbę razy. Pozostałe
składniki mają moduł jeden i najwyżej $m$ różnych
indeksów. Ich liczba nie przekracza
$(2m)^{2m}N^m$ (wybór wzorca indeksów i samych indeksów).
Nierówność Höldera na przestrzeni probabilistycznej daje
zatem
\[
 \mathbb E\left|\sum_j\varepsilon_j e^{ijx_1}\right|^p
 \le (2m)^p N^{p/2}.
\]
Całkując i stosując Tonellego, znajdujemy wybór znaków
taki, że $\|f_N\|_p\le 2m\sqrt N\|\phi\|_p$.
Gdyby rozważane oszacowanie zachodziło, to
$N^{1/q-1/2}$ byłoby ograniczone niezależnie od $N$,
co jest niemożliwe, gdy $q<2$.

Dla $p=\infty$ korzystamy z jawnej konstrukcji
wielomianów Rudina--Shapira:
$P_0=Q_0=1$ oraz
\[
 P_{n+1}(z)=P_n(z)+z^{2^n}Q_n(z),\qquad
 Q_{n+1}(z)=P_n(z)-z^{2^n}Q_n(z).
\]
Indukcyjnie oba mają po $2^n$ współczynników
równych $1$ lub $-1$, ponieważ nośniki potęg
w obu składnikach są rozłączne. Dla $|z|=1$ mamy
$|P_{n+1}|^2+|Q_{n+1}|^2=2(|P_n|^2+|Q_n|^2)$,
więc $|P_n(z)|\le\sqrt{2^{n+1}}$.
Biorąc $N=2^n$ i $f_N(x)=\phi(x)e^{ix_1}P_n(e^{ix_1})$,
otrzymujemy $\|f_N\|_\infty\le\sqrt{2N}\|\phi\|_\infty$,
podczas gdy $\|\widehat f_N\|_1=N\|\psi\|_1$.
To również wyklucza istnienie stałej $C$.
\end{proof}
